\documentclass{article}

\usepackage[spanish]{babel}
\usepackage{amsmath,amssymb}

\newcommand{\N}{\mathbb{N}}
\newcommand{\Z}{\mathbb{Z}}
\newcommand{\Q}{\mathbb{Q}}
\newcommand{\R}{\mathbb{R}}
\newcommand{\C}{\mathbb{C}}

\newtheorem{theorem}{Teorema}

\title{Estadistica y Probabilidad - Variable Aleatoria}
\author{Benjamín Rivas Beltrán}
\date{\today}

\begin{document}

\maketitle

\section{Variable Aleatoria y Distribución de Probabilidad}

Una variable aleatoria (v.a.) es una función que toma cada resultado de un experimento aleatorio $x_i$ y le asigna un valor real.

\begin{table}[h]
\centering
\begin{tabular}{|c|c|c|c|c|c|}
\hline
$x_i$ & $x_1$ & $x_2$ & $\ldots$ & $x_n$ & $\sum$ \\
\hline
$p(x_i)$ & $p(x_1)$ & $p(x_2)$ & $\ldots$ & $p(x_n)$ & 1 \\
\hline
\end{tabular}
\end{table}

Se llama \textbf{distribución de probabilidad} a la función $p(x_i)$ que asigna a cada valor $x_i$ la probabilidad de que la variable aleatoria tome ese valor.

\subsection{Propiedades de la distribución de probabilidad}
\begin{itemize}
  \item $0 \leq p(x_i) \leq 1$ para todo $i$.
    \item $\sum_{i=1}^{n} p(x_i) = 1$.
\end{itemize}

\subsection{Función de distribución}
Si tenemos una distribución de probabilidad, podemos determinar su función de distribución $F(x)$, tal que $F(x_i) = P(X \leq x_i)$, es decir, la probabilidad de que la variable aleatoria tome un valor menor o igual a $x_i$.

Entonces, se tiene que:
\begin{align*}
F(x_1) &= P(X \leq x_1) = p(x_1) \\
F(x_2) &= P(X \leq x_2) = p(x_1) + p(x_2) \\
&\vdots \\
F(x_i) &= P(X \leq x_i) = \sum_{j=1}^{i} p(x_j) \\
&\vdots \\
F(x_n) &= P(X \leq x_n) = \sum_{j=1}^{n} p(x_j) = 1
\end{align*}

Podemos calcular probabilidades con la distribución de probabilidad o con la función de distribución, y en tal  caso, se tiene que:
\begin{align*}
P(a < x < b) &= F(b) - F(a) - p(x = b) \\
P(a < x \leq b) &= F(b) - F(a) \\
P(a \leq x < b) &= F(b) - F(a) + p(x = a) - p(x = b) \\
P(a \leq x \leq b) &= F(b) - F(a) + p(x = a)
\end{align*}

\subsection{Valor esperado, varianza y función generadora de momentos}

Para cualquier distribución de probabilidad, podemos determinar su valor esperado (promedio), y se denota por $E[X]$

\begin{align*}
E[X] &= \sum_{i=1}^{n} x_i p(x_i)
\end{align*}

También podemos determinar su varianza y desviación estándar, y se denotan por $V(X)$ y $\sigma_X$ respectivamente.

\begin{align*}
V(X) &= E[X^2] - (E[X])^2 \\
\sigma_X &= \sqrt{V(X)}
\end{align*}

Donde a $E[X^2]$ se le llama el segundo momento de la variable aleatoria, y se calcula de la siguiente manera:
\begin{align*}
E[X^2] &= \sum_{i=1}^{n} x_i^2 p(x_i)
\end{align*}

En general, el $k$-ésimo momento de la variable aleatoria se denota por $E[X^k]$ y se calcula de la siguiente manera:
\begin{align*}
E[X^k] &= \sum_{i=1}^{n} x_i^k p(x_i)
\end{align*}

Por último, podemos determinar la función generadora de momentos $M_X(t) = E[e^{tX}]$, y se calcula de la siguiente manera:
\begin{align*}
M_X(t) &= E[e^{tx_1}] = \sum_{i=1}^{n} e^{tx_i} p(x_i)
\end{align*}

\begin{theorem}
Si tenemos un función generadora de momentos $M_X(t)$, su $k$-ésima derivada evaluada en $t=0$ nos da el $k$-ésimo momento de la variable aleatoria, es decir:
\begin{align*}
\frac{d^k}{dt^k} M_X(t) \bigg|_{t=0} = E[X^k]
\end{align*}
\end{theorem}

\subsection{Propiedades del valor esperado}

Sea $x$ una variable aleatoria y $C, a, b \in \R$ constantes, entonces se cumplen las siguientes propiedades:

\begin{itemize}
    \item $E[C] = C$.
    \item $E[Cx] = C E[X]$.
    \item Si tenemos una combinación lineal $y = a + b x$, entonces:
      \begin{align*}
      E[y] &= E[a + b x] \\
      &= E[a] + E[b x] \\
      &= a + b E[X]
      \end{align*}
\end{itemize}

\subsection{Propiedades de la varianza}
Sea $X$ una variable aleatoria y $C, a, b \in \R$ constantes, entonces se cumplen las siguientes propiedades:

\begin{itemize}
    \item $V(C) = 0$.
    \item $V(CX) = C^2 V(X)$.
    \item Si tenemos una combinación lineal $Y = a + bX $, entonces:
      \begin{align*}
      V(Y) &= E[Y^2] - (E[Y])^2 \\
      &= E[(a + bX)^2] - (E[a + bX])^2 \\
      &= E[a^2 + 2abX + b^2 X^2] - (a + b E[X])^2 \\
      &= a^2 + 2ab E[X] + b^2 E[X^2] - (a^2 + 2ab E[X] + b^2 (E[X])^2) \\
      &= a^2 + 2ab E[X] + b^2 E[X^2] - a^2 - 2ab E[X] - b^2 (E[X])^2 \\
      &= b^2 (E[X^2] - (E[X])^2) \\
      &= b^2 V(X)
      \end{align*}
\end{itemize}

\end{document}